\pg{343} 66 O ente, que é acto puro, é infinito; o ente, que é \textbf{infinito}, é acto puro.

a) \emph{O ente}, que é \emph{acto puro}, \emph{é infinito}. --- O ente, que é \emph{acto puro}, tem o ser, n\&o recebido na essencia, mas por si subsistente. Ora o ente, que tem o ser, não recebido na essencia, mas por si subsistente, é \emph{infinito}; pois a limitação do ser deriva da es-

\emph{H} sencia, em que é recebido. Logo o ente, que é \emph{acto puro}, \emph{é infinito}.

b) \emph{O ente}, \emph{que é infinito}, \emph{é acto puro}. --- O ente \emph{infinito} é dotado de toda a perfeição possivel, é a propria perfeição. Ora o ente, dotado de toda a perfeição possivel, é necessariamente \emph{acto puro}; porque só se lhe poderia aliribuir alguma potencialidade, se Ihe faltasse alguma perfeição, mas então deixaria de ser infinito. Logo o ente, que é \emph{infinito}, é \emph{acto puro} (\textasciicircum{}).

67. O ente, que \textbf{consta} de acto e de potencia, é finito; o ente. que é finito, \textbf{consta} de acto e de potencia,

a) \emph{O ente}, \emph{que-consta-de-acto-e-de-potencia}, \emph{é finito}. --- --- O. ente; A

que consta de \emph{acto} e.de \emph{potencia}, tem o proprio ser recebido na essencia. Ora o ser recebido na essencia é limitado pela propria essencia, em que é recebido. Logo o ente, que consta de \emph{acto} e de \emph{potencia}, \emph{é finito} \footnote{A verdade d'estas duas proposições apparece evidente pela simples declaração dos termos. \emph{Acto} € perfeição; acto \emph{puro} é perfeição pura, sem sombra de imperfeição, e por isso illimitada, infinita. Vice-versa, o ente infinito. que é o proprio ser, é \emph{acto puro}; porque, se fosse composto de potencia e de alo, seria limitado, deixaria de ser infinito. --- Possuindo todas as perfeições possiveis, o infinito possue todas as perfeições dos entes finitos, tanto as \emph{simples} como as \emph{snistas}, --- aquellas de um modo \emph{formal} e \emph{eminente}, estas de um modo \emph{virtual}! e \emph{eminente}. Por isso, os termos, que exprimem as perfeicoes \emph{séxples}, quando atiribuidos a Deus, não egualam nem limitam a coisa significada, como quando se attribuem ás creaturas, mas deixam-na incomprehensivel e superior á significação dos proprios lermos. (\emph{Sum}. \emph{Th}., p. I, q. 13, a. 5).}.\footnote{Só o \emph{acto puro}, sendo o proprio ser subsistente, € infinito. Onde se \emph{encontra} a potencialidade, ahi se encontra a limitação. --- Alem d'isso, o \emph{infinito} não é o resultado da addic\&o de partes. Portanto, onde se encontra \& parte, ahi se encontra a limitação. Nem vale recorrer a um numero immenso de partes ; o finito, addicionado ao finito, dá sempre \emph{o} finito.\par D'ahi os seguintes corollarios da mais alta importancia\par 1.º) \emph{Não pode haver} um \emph{infinito} maior \emph{do} que outro \emph{infinito}. --- O infinito não admitte o mais e o menos; e por isso, se fosse \emph{excedido} por outro infnito\pgn{344}, deixaria de ser tal, E não só não admitte um outro infinito maior, mas nem comporta outro \emph{egual}; porque, se fosse duplice ou multiplice, deixarja de ser infinito. Por isso, deve ser unico. - Quando se diz que não pode haver um infinito maior que*outro infinito, falla-se no infinito, que é tal em acto ou absolutamente, e não se falla no infinito, que só é tal em potencia, emquanto é sempre capaz de receber novas partes, ou novas unidades, e que em \emph{acto}, ou na \emph{renlidode}, é sempre finito. --- O infinito \emph{em potencia} não é infinito sob todos os aspectos nem em todas as ordens, mas é tal “apenas sob um certo aspecto e numa só ordem; nem é uma substancia, mas é um attributo accidental das coisas, que se dizem infinitas, e por isso pode multiplicar-se, conforme se multiplica o sujeito, a que se attribue. Assim, podemos imaginar muitas linhas infinitas no comprimento, infinitas series de numeros pares e impares, infinitas series de effeitos possiveis em cada uma das causas, infinitos individuos possiveis em cada especie, etc. --- Como tambem, este infinito em \emph{potencia} não comporta outro infinito maior na mesma \emph{ordem}, em que é infinito; mas, por causa de um sujeito maior, pode admittir outro infinito, maior do que elle, numa \emph{outra ordem}. Assim, as especies dos numeros: pares são infinitas, e lambem infinitas são as especies dos numeros impares; e, todavia, a colleccko dos numeros pares e impares é maior do que a dos numeros pares. .Bsta theoria é de S. Thomaz (\emph{Quodlib}. 1X, q. 1, a. 1 ad 1).\par 2,0) \emph{Nenhuma} creatura, \emph{pode ser infinita}. quanto á essencia, --- O ente, que é infinito quanto é \emph{essencia}, é o proprio Ser subsistente, acto purissimo, perfeição sem limite, Ora nenhuma creatura, sendo contingente, é o proprio ser subsistente, o acto purissimo, a propria perfeição. Logo nenhuma creatura pode ser infinita \emph{guanto} d essencia. Em menos palavras : creatura infinita na \emph{essencia} é uma contradicgio. Por isso S. Thomaz nega que Deus possa' produzir uma creatura infinita quanto á essencia. (Sum. \emph{Th}., p. I, q. 4, a. 2 ad 1). --- Para a refutação das objecQóes, basta considerar --- a) que um effeito infinito é um effeito necessario. independente de todo e qualquer ente, e por isso é um effeito que não é effeito; --- 5) que o effeito da causa infinita deve demonstrar e demonstra a infinidade da causa, --- não pela infinidade da propria essencia (o que é absurdo). --- mas pela infinidade da força necessaria para a sua producção, porque só uma força infinita \emph{qoos} crear, ou tirar do nada as coisas.\par 3.º) \emph{Nenhuma creatura pode ser} a mais \emph{perfeita de todas}. --- \emph{Durando}, \emph{Henrique de}. \emph{Gand}, \emph{Gabriel Bie}, Occam, \emph{Arriaga}, e muitos \emph{Scotisías e Nominalisios} ensinaram que pode haver uma creatura tão perfeita, que não possa baver outra mais perfeita. \emph{S}. \emph{Thomas} e, geralmente, os \emph{Xscholasticos} são contrarios a essa doutrina, e com razão. Na verdade, uma creatura, porquanto se suppouha perfeita, nunca pode exgottar a participabilidade da perfeição divina, e entre essa creatura e Deus existe sempre uma distancia infinita, de modo que podem dar-se ainda infinitas outras creaturas, cada uma das quaes seja mais perfeita do que a que se diz perfeitissima, sem que nenhuma chegue a exgottar a divina partieipabilidade. Portanto, uma creatura, que seja a mais perfeita de todas, está fóra da serie dos possiveis. ---\pgn{345}Egualmente repugna uma creatura que seja a mais imperfeita de todas, de modo que não se conceba outra mais imperfeita. Toda. 2. creatura, embora constituida numa especie imperfeitissima, comtudo dista infinitamente do nada, e por isso, abaixo d' ella, pode haver outras creaturas sem fim.\par 4.º) \emph{Nenhuma crentura pode ser dotada de força} ou actividade \emph{infinita}. --- É um corollario da primeira proposição. A força ou a actividade deriva da essencia, e é proporcionada a ella. Se a essencia das creaturas é necessariamente finita, tambem a sua força ou actividade é finita. (Cf. C. Gent., 1, 43), --- Por isso, a creatura racional. embora tenda, como para o seu fim, para o Bem infinito, comtudo só o alcança por um modo finito, isto é, por actos finitos de faculdades finitas, que são a intelligencia e a vontade,\par 5.º) \emph{Uma grandeza} OU \emph{extensão} acinalmente \emph{infinita repugna}. --- Uma essencia determinada é dotada de propriedades determinadas. Ora a grandeza ou extensão é uma propriedade do corpo, e todo o corpo tem uma essencia determinada. Logo a grandeza ou extensão do corpo é determinada, e por isso finita. D' onde se segue que a extensão de um corpo tem certos limites, que não é possivel transpor, sem que o proprio corpo seja physica-:. mente destruido, embora a extensão, considerada em si, seja ulteriormente divisivel. --- O que dizemos da \emph{quantidade} pode applicar-se \& \emph{qualidade}: Repugna que uma qualidade seja Infinita na sua intensidade; pois, como dissemos, uma propriedade deve ser e é sempre proporcionada ao sujeito, de que deriva e em que reside --- Todavia, a grandeza, considerada de um modo abstracto, diz-se \emph{potencialmente} infinita, emquanto pode augmentar sem fim,\par 6.º) \emph{Uma multidão actualmente infinita repugna}. --- Toda a \emph{multidão} é uma especie de numero, resulta da addição de unidades, e é sempre capaz de augmento é de diminuição. Ora o \emph{infinito} --- não pode \emph{ser} uma especie de numero (pois O numero é uma multidão medida: pela unidade, emquanto o infinito, por ser tal, não tem limites e por isso não pode ser medido), --- não resulta da addição de unidades (porque as unidades são partes finitas, e estas não podem dar o infinito), --- não é capaz de augmento e de diminuição (porque uma multidão, a qual pode accrescentar-se uma unidade, já se presuppõe finita; assim como uma multidão, que se suppõe infinita, se lhe se podesse tirar uma unidade, deixaria de ser infinita). Logo uma multidão actualmente infinita repugna.\par Dissemos --- \emph{actualmente} infinita, para se denotar que a multidão infinita repugna, quando todos os seus elementos se consideram \emph{em acto}, isto é, na ordem \emph{da existencia real}; porque, se os elementos se consideram em potencia, isto é, na ordem \emph{da possibilidade}, neste caso a multidão infinita não .repugna. Uma multidão infinita de entes \emph{possiveis} é possivel. Na verdade, o augmento de uma multidão está na razão directa da divisão da grandeza ou extensão; visto que, quanto mais se divide uma grandeza, tanto mais cresce o numero das partes. Ora, assim como a grandeza é divisivel sem fim, assim tambem a multidão cresce sem fim. isto é, a infinidade \emph{potencial} da multidão é uma consequencia da infipidade \emph{potencial} da grandeza ou extensão. Logo uma multidão infinita \emph{em potencia} é possivel. (\emph{Sum}. \emph{Th}., q. 1. 9. 7, a. 4). --- Agora uma advertencia. A multidão, que se diz infinita \emph{em potencia}, emquanto pode ser dividida em partes, cada vez mais pequenas, \emph{ou}\pgn{346}pode receber outras unidades, sem nunca chegar ao fim, é, de facto, sempre limitada, porque constituída por um. numero finito de unidades. Por isso \emph{o} infinito \emph{em potencia} costuma chamar-se, mais propriamente, \emph{indefinito}.}

\pg{344} 544

b) O \emph{ente}, que é \emph{finito}, consta \emph{de acto} e \emph{de potencia}. --- O ente, que é \emph{finito}, tem o proprio ser limitado, Ora o ser emtanto é \pg{345}

N 